\documentclass[tikz,border=8pt]{standalone}
\usepackage{tikz}
\usepackage{amsmath}
\usepackage{amssymb}
\usepackage{pgfplots}
\pgfplotsset{compat=1.18}
\usepackage{ctex}
\usetikzlibrary{calc,positioning,arrows.meta,matrix}

% LC 56 Merge Intervals
% 原始区间 [1,3],[2,6],[8,10],[15,18] → 合并后 [1,6],[8,10],[15,18]

\begin{document}
\begin{tikzpicture}[
  every node/.style={font=\small},
]

%% ── 标题 ──────────────────────────────────────────────────
\node[font=\large\bfseries] at (8.5, 1.3)
  {区间合并（LC 56 Merge Intervals）};
\node[font=\small,gray] at (8.5, 0.65)
  {先按左端点排序，再扫描合并：当前区间左端 $\leq$ 前一区间右端则合并};

%% ────────────────────────────────────────────────────────
%% 坐标说明：轴起点 ox=0.8，scale=0.9（每单位=0.9cm）
%% 数值 v 的 x = 0.8 + v*0.9
%% 区间 [a,b] 的 xL = 0.8+a*0.9, xR = 0.8+b*0.9, 中点 = 0.8+(a+b)/2*0.9
%% ────────────────────────────────────────────────────────

%% ── 上方区域（原始区间）─────────────────────────────────
\node[font=\scriptsize\bfseries,gray,anchor=east] at (0.6, -0.5) {原始：};

% 数轴（y=-0.5）
\draw[->,>=Stealth,gray!50,thick] (0.8, -0.5) -- (17.1, -0.5);
% 刻度
\foreach \v in {0,1,2,3,4,5,6,7,8,9,10,11,12,13,14,15,16,17,18} {
  \pgfmathsetmacro\tx{0.8+\v*0.9}
  \draw[gray!40,thin] (\tx, -0.57) -- (\tx, -0.43);
}
% 数字（每3个）
\foreach \v in {0,3,6,9,12,15,18} {
  \pgfmathsetmacro\tx{0.8+\v*0.9}
  \node[font=\tiny,gray] at (\tx, -0.82) {\v};
}

% 区间 [1,3]，row y=-1.5
% xL=0.8+1*0.9=1.7, xR=0.8+3*0.9=3.5, mid=0.8+2*0.9=2.6
\draw[draw=blue!60,fill=blue!10,rounded corners=1pt,thick]
  (1.7, -1.68) rectangle (3.5, -1.32);
\node[font=\scriptsize\bfseries,blue!70] at (2.6, -1.5) {[1,3]};

% 区间 [2,6]，row y=-2.4
% xL=0.8+2*0.9=2.6, xR=0.8+6*0.9=6.2, mid=0.8+4*0.9=4.4
\draw[draw=blue!60,fill=blue!10,rounded corners=1pt,thick]
  (2.6, -2.58) rectangle (6.2, -2.22);
\node[font=\scriptsize\bfseries,blue!70] at (4.4, -2.4) {[2,6]};

% 区间 [8,10]，row y=-1.5
% xL=0.8+8*0.9=8.0, xR=0.8+10*0.9=9.8, mid=0.8+9*0.9=8.9
\draw[draw=gray!60,fill=gray!10,rounded corners=1pt,thick]
  (8.0, -1.68) rectangle (9.8, -1.32);
\node[font=\scriptsize\bfseries,gray!70] at (8.9, -1.5) {[8,10]};

% 区间 [15,18]，row y=-1.5
% xL=0.8+15*0.9=14.3, xR=0.8+18*0.9=17.0, mid=0.8+16.5*0.9=15.65
\draw[draw=gray!60,fill=gray!10,rounded corners=1pt,thick]
  (14.3, -1.68) rectangle (17.0, -1.32);
\node[font=\scriptsize\bfseries,gray!70] at (15.65, -1.5) {[15,18]};

% 重叠标注：[1,3] 和 [2,6] 的重叠区间 x=2.6..3.5
\draw[red!50,dashed,thin] (2.6, -1.32) -- (2.6, -2.22);
\draw[red!50,dashed,thin] (3.5, -1.32) -- (3.5, -2.22);
\node[font=\tiny,red!70,anchor=west] at (3.6, -1.9) {重叠！};

% 说明框（右侧）
\node[draw=gray!40,fill=white,thin,rounded corners=2pt,
      font=\scriptsize,inner sep=4pt,align=left]
  at (13.5, -2.0)
  {[1,3]与[2,6]：\\左端 $2 \leq 3$，合并\\$\to [1,\max(3,6)]=[1,6]$};

%% ── 中间箭头 ─────────────────────────────────────────────
\draw[->,>=Stealth,gray!60,thick] (8.5, -3.1) -- (8.5, -4.1)
  node[midway,right,font=\scriptsize,gray] {合并后};

%% ── 下方区域（合并后区间）───────────────────────────────
\node[font=\scriptsize\bfseries,gray,anchor=east] at (0.6, -4.8) {合并后：};

% 数轴（y=-4.8）
\draw[->,>=Stealth,gray!50,thick] (0.8, -4.8) -- (17.1, -4.8);
\foreach \v in {0,1,2,3,4,5,6,7,8,9,10,11,12,13,14,15,16,17,18} {
  \pgfmathsetmacro\tx{0.8+\v*0.9}
  \draw[gray!40,thin] (\tx, -4.87) -- (\tx, -4.73);
}
\foreach \v in {0,3,6,9,12,15,18} {
  \pgfmathsetmacro\tx{0.8+\v*0.9}
  \node[font=\tiny,gray] at (\tx, -5.12) {\v};
}

% 合并后区间 [1,6]
% xL=1.7, xR=6.2, mid=3.95
\draw[draw=red!60,fill=red!10,rounded corners=1pt,very thick]
  (1.7, -5.02) rectangle (6.2, -4.58);
\node[font=\scriptsize\bfseries,red!70] at (3.95, -4.8) {[1, 6]};

% 合并后区间 [8,10]
% xL=8.0, xR=9.8, mid=8.9
\draw[draw=gray!70,fill=gray!10,rounded corners=1pt,thick]
  (8.0, -5.02) rectangle (9.8, -4.58);
\node[font=\scriptsize\bfseries,gray!80] at (8.9, -4.8) {[8,10]};

% 合并后区间 [15,18]
% xL=14.3, xR=17.0, mid=15.65
\draw[draw=gray!70,fill=gray!10,rounded corners=1pt,thick]
  (14.3, -5.02) rectangle (17.0, -4.58);
\node[font=\scriptsize\bfseries,gray!80] at (15.65, -4.8) {[15,18]};

%% ── 红色虚线：来源追踪（[1,3]+[2,6] → [1,6]）───────────
% 从上方 xL=1.7 连到下方 xL=1.7
\draw[red!40,dashed,thin] (1.7, -2.58) -- (1.7, -4.58);
% 从上方 xR=6.2(即[2,6]的右端) 连到下方 xR=6.2
\draw[red!40,dashed,thin] (6.2, -2.58) -- (6.2, -4.58);

%% ── 分隔线 ──────────────────────────────────────────────
\draw[gray!25,dashed] (0.0, -6.0) -- (17.5, -6.0);

%% ── 算法说明 ─────────────────────────────────────────────
\node[draw=gray!50,fill=white,rounded corners=3pt,font=\small,
      inner sep=8pt,align=left]
  at (8.5, -7.1)
  {\textbf{Merge Intervals 算法（时间 $O(n\log n)$）：}\\[4pt]
   1. 按左端点排序区间列表\\[2pt]
   2. 用结果列表 \texttt{res}，初始放入第一个区间\\[2pt]
   3. 遍历余下区间 \texttt{[l, r]}：\\
   \quad 若 $l \leq \texttt{res[-1][1]}$（有重叠）$\to$ 更新右端 \texttt{res[-1][1] = max(res[-1][1], r)}\\
   \quad 否则（无重叠）$\to$ 直接追加 \texttt{[l, r]} 到 \texttt{res}};

\end{tikzpicture}
\end{document}
